\documentclass{article}
\usepackage[utf8]{inputenc}
\usepackage{geometry}
\geometry{
	a4paper,
	total={170mm,257mm},
	left=20mm,
	top=20mm,
}
\usepackage{graphicx}
\usepackage{titling}
\usepackage[american,RPvoltages]{circuiTikz}
\usepackage{tikz, tikz-cd}
\usepackage{pgfplots}
\usepackage{siunitx}
\usepackage{float}
\usepackage{enumitem}
\usepackage{amsmath,amsfonts,stmaryrd,amssymb} % Math packages
\title{Diodos rectificadores}
\author{Rodrigo}
\date{July 2024}

\usepackage{fancyhdr}
\fancypagestyle{plain}{%  the preset of fancyhdr 
	\fancyhf{} % clear all header and footer fields
	\fancyfoot[L]{\thedate}
	\fancyhead[L]{Problem Resolution}
	\fancyhead[R]{\theauthor}
}
\makeatletter
\def\@maketitle{%
	\newpage
	\null
	\vskip 1em%
	\begin{center}%
		\let \footnote \thanks
		{\LARGE \@title \par}%
		\vskip 1em%
		%{\large \@date}%
	\end{center}%
	\par
	\vskip 1em}
\makeatother

\usepackage{lipsum}  
\usepackage{cmbright}

\begin{document}
	
	\maketitle
	
	\noindent\begin{tabular}{@{}ll}
		Author & \theauthor\\
		Chapter &  Basic Circuit Analysis\\
		Problem & 2.68
	\end{tabular}
	
	\section*{Resumen}
	Los diodos son extremadamente usados en rectificadores, la mayoría de los equipos electrónicos y circuitos requieren de fuentes de DC para que operen correctamente. Las baterías recargables pueden ser usadas en estas aplicaciones, sin embargo, sólo ofrecerían una potencia limitada y un voltaje inestable conforme se vayan descargando las celdas. Las fuentes de alimentación de DC más comunes se realizan a partir de rectificadores AC/DC.
	Un Rectificador es un circuito que convierte una señal de AC en una señal unidireccional. Los rectificadores AC/DC sin controlar usualmente consisten de circuitos con diodos los cuales pueden ser clasificados en los siguientes grupos.
	
	\begin{itemize}[noitemsep]
		\item Rectificadores de media onda
		\item Rectificadores de onda completa
		\item Rectificadores de tres fases
		\item Rectificadores multipulso
		\item Rectificadores de Corrección de Factor de Potencia (PFC)
		\item Rectificadores tipo boost modulados por ancho de pulso 
	\end{itemize}

	
	\section*{Introducción}
	La entrada de un diodo rectificador es un voltaje de AC, el cual puede ser un voltaje monofásico o trifásico, usualmente de una onda senoidal pura. Un voltaje de entrada monofásico $v(t)$ puede ser expresado como:
	
	\begin{align}
		v(t) = \sqrt{2}V \sin{\omega t} = V_m \sin{\omega t}
	\end{align}
	
	\noindent	Donde: \\
	
	\noindent\begin{tabular}{@{}ll}
		$v(t)$ & Es el voltaje de entrada instantáneo\\
		V &  Es el valor rms\\
		$V_m$ & Es la amplitud\\
		$\omega$ & Es la frecuencia angular donde $\omega= 2\pi f$ ($f$ es la frecuencia de la fuente)
	\end{tabular}
	
	\section*{Rectificadores de media onda}
	Un circuito rectificador de media onda consiste de un voltaje de entrada monofásico de AC y un diodo. El circuito se muestra en la figura 1.
	
	\begin{figure}[H]
	\centering	
	
	\begin{tikzpicture}[scale=0.8]
		\draw (0,0) to[sV, l=${v=\sqrt{2}V \sin{\omega t} }$, i=$i$] ++(0,5)
			  to[ccgsw] ++(3,0)
			  to[D*, v=$v_D$, l=$D$] ++(3,0)
			  to[R, v=$v_o$, l=$R$] ++(0,-5)
			  to[short] (0,0); 
			  
		\pgfplotsset{width=9cm,compat=1.18}
		\begin{axis}[xshift=9cm,
			 yshift=-0.5,
			 clip=false,
			 axis y line=left,
			 axis x line=middle, 
			 xmin=0,xmax=2.2*pi,
			 ymin=-1.5,ymax=1.5,
			 xlabel=$\omega t$,
			 xtick={3.14,6.28},
			 xticklabels={$\pi\,$, $\,\,\,2\pi$},
			 ytick={1},
			 yticklabels={$\sqrt{2}V$},
			 ]
			\addplot[smooth, blue, mark=none, domain=0:2*pi, samples=200] {sin(deg(x))};
		\end{axis}
		
		\draw (9.5,6) node[above right, xshift=-20pt]{$v$};
	\end{tikzpicture}
	\caption{}	
	\label{fig:figura1}
	\end{figure}
	
	\begin{figure}[H]
	\centering	
	
	\begin{tikzpicture}[scale=0.8]
			\pgfplotsset{width=9cm,compat=1.18}
		\begin{axis}[
			clip=false,
			axis y line=left,
			axis x line=middle, 
			xmin=0,xmax=2*pi,
			ymin=-1.5,ymax=1.5,
			xlabel=$\omega t$,
			xtick={3.14,6.28},
			xticklabels={$\pi\,$, $\,\,\,2\pi$},
			ytick={1},
			yticklabels={$\sqrt{2}V$},
			]
			\addplot[smooth, blue, mark=none, domain=0:pi, samples=200] {sin(deg(x))};
		\end{axis}
		
		\draw (0.3,6) node[above right, xshift=-20pt]{$v_o$};
		
		\pgfplotsset{width=9cm,compat=1.18}
		\begin{axis}[xshift=9cm,
			yshift=-0.5,
			clip=false,
			axis y line=left,
			axis x line=middle, 
			xmin=0,xmax=2*pi,
			ymin=-1.5,ymax=1.5,
			xlabel=$\omega t$,
			xtick={3.14,6.28},
			xticklabels={$\pi\,$, $\,\,\,2\pi$},
			ytick={1},
			yticklabels={$\frac{\sqrt{2}V}{R}$},
			]
			\addplot[smooth, blue, mark=none, domain=0:pi, samples=200] {sin(deg(x))};
		\end{axis}
		
		\draw (9.5,6) node[above right, xshift=-20pt]{$i$};
	\end{tikzpicture}
	\caption{}	
	\label{fig:figura2}
\end{figure}

	\begin{figure}[H]
	\centering	
	
	\begin{tikzpicture}[scale=0.8]
		\draw (0,0) to[sV, l=${v=\sqrt{2}V \sin{\omega t} }$, i=$i$, v=$$] ++(0,5)
		to[ccsw] ++(5,0)
		to[R, v=$v_O$, l=$R$] ++(0,-5) coordinate(p1)
		to[short] (0,0); 
		\draw (2,-1) node[above right]{$v > 0$};
		\draw (p1) ++(3,0) coordinate(p2) to[sV] ++(0,5)
		to[cosw] ++(5,0)
		to[R, v=$v_O$, l=$R$] ++(0,-5)
		to[short] ++(-5,0);
		\draw (10,-1) node[above right]{$v < 0$};
	\end{tikzpicture}
	\caption{}	
	\label{fig:figura3}
\end{figure}

El diagrama del circuito del rectificador de media onda se muestra en la figura 1(a). Consideremos el intervalo 1 como $0 \leq \omega t \leq \pi $ durante el semiciclo positivo del voltaje de entrada. El diodo $D$ conduce y se comporta como un cortocircuito, como se muestra en la figura 3(a). El voltaje de entrada aparece a través de la resistencia $R$, esto es, el voltaje de salida sería:

\begin{align}
	v_O = V_m \sin{\omega t}, \quad 0 \leq \omega t \leq \pi
\end{align}

Si incluimos la caída de voltaje del diodo ($V_D = \qty{0.7}{\volt}$) el voltaje pico de salida $V_m$ sería reducido a ($V_m - V_D$) y el voltaje instantáneo de salida sería:

\begin{align}
	v_O = (V_m - V_D) \sin{\omega t}, \quad 0 \leq \omega t \leq \pi
\end{align}

Consideremos $\pi \leq \omega t \leq 2 \pi$ como el intervalo 2 durante el semiciclo negativo del voltaje de entrada. El diodo $D$ se encontrará polarizado en inversa, por lo que se comportará como un circuito abierto, como se muestra en la figura 3(b). Por lo tanto el voltaje de salida $v_O$ sería igual a cero, esto es:

\begin{align}
	v_O = 0, \quad \pi \leq \omega t \leq 2 \pi
\end{align}

El voltaje de salida promedio $V_{O(av)}$ se encuentar usando la siguiente ecuación:

\begin{align}
	V_{O(av)} = \frac{1}{T} \int_{0}^{T} f(t) d(\omega t)
\end{align}

Por lo tanto, a partir de la gráfica de la figura 2(a) y de las ecuaciones (2) y (4) podemos definir el voltaje de salida como

\begin{align}
	v_O= \left\{ 
	\begin{array}{lc} V_m \sin{\omega t}, & 0 \leq \omega t \leq \pi \\
	0, & \pi \leq \omega t \leq 2\pi
	 \end{array} \right.
\end{align}


Así

\begin{align}
	V_{O(av)} &= \frac{1}{2\pi} \int_{0}^{\pi} V_m \sin{\omega t} \,\, d(\omega t) \nonumber \\
	 &= \frac{V_m}{2\pi} \int_{0}^{\pi} \sin{\omega t} \,\, d(\omega t) \nonumber \\
	&= - \left.{\left[ \frac{V_m}{2\pi} \cos({\omega t}) \right]}\right|_{\;0}^{\;\pi} \nonumber \\
	&= - \frac{V_m}{2\pi} \left[ cos (\pi) - \cos(0) \right] \nonumber \\
	&= - \frac{V_m}{2\pi} (-1 - 1) \nonumber \\
	&= \frac{2 V_m}{2\pi} \nonumber
\end{align}

O bien

\begin{align}
	V_{O(av)} = \frac{V_m}{\pi} \approx 0.318V_m
\end{align}

Por lo tanto, la corriente promedio $I_{O(av)}$ en la resistencia de carga $R$ puede encontrarse mediante la ley de Ohm

\begin{align}
	I_{O(av)} = \frac{V_{O(av)}}{R} = \frac{V_m}{\pi R} = \frac{0.318V_m}{R}
\end{align}

El voltaje rms de una señal está dado por:

\begin{align}
	V_{rms} = \sqrt{\frac{1}{T} \int_{0}^{T} {f^2 (t) dt}}
\end{align}

Así, el voltaje rms de salida del rectificador de media onda puede escribirse como

\begin{align}
	V_{o(rms)} &= \sqrt{\frac{1}{2\pi} \int_{0}^{\pi} v_O^2 \, d(\omega t)} = \sqrt{\frac{1}{2\pi} \int_{0}^{\pi} V_m^2 \sin^2(\omega t) \, d(\omega t)} \nonumber \\
	&= \sqrt{\frac{V_m^2}{2\pi} \int_{0}^{\pi} \left[\frac{1}{2} - \frac{1}{2} \cos(2 \omega t)\right] \, d(\omega t)} \nonumber
\end{align}

Donde

\begin{align}
	\int_{0}^{\pi} \left[\frac{1}{2} - \frac{1}{2} \cos (2\omega t) d(\omega t) \right] &= \int_{0}^{\pi} \frac{1}{2} d(\omega t) - \int_{0}^{\pi} \frac{1}{2} \cos (2\omega t) d(\omega t) \nonumber \\
	&= \left.{\left[ \frac{\omega t}{2} \right]}\right|_{\;0}^{\;\pi} - \left.{\left[ \frac{1}{4} \sin (2\omega t) \right]}\right|_{\;0}^{\;\pi} \nonumber \\
	&= \frac{\pi}{2} - \frac{1}{4} \sin (2\pi) + \frac{1}{4} \sin 0 = \frac{\pi}{2} \nonumber
\end{align}

Por último
\begin{align}
	V_{o(rms)} &= \sqrt{ \frac{V_m^2}{2\pi} \left( \frac{\pi}{2} \right) } = \sqrt{ \frac{V_m^2}{4}  } \nonumber
\end{align}

\begin{align}
	V_{o(rms)} = \frac{V_m}{2} = 0.5V_m
\end{align}

Y la corriente rms de la carga $I_{o(rms)}$ está dada por

\begin{align}
	I_{o(rms)} = \frac{V_{o(rms)}}{R} = \frac{V_m}{2R} = \frac{0.5V_m}{R}
\end{align}

\subsection*{Rectificador de media onda con el diodo real}
Asumiendo una aproximación lineal para el diodo en dirección directa, el diodo comienza a conducir en el medio ciclo positivo, cuando $V_m \sin \omega t = V_{D0}$, mientras el diodo conduce, la corriente en la carga será igual a:

\begin{align}
	i_L = \frac{V_m \sin \omega t - V_{D0}}{R_L + r_D}
\end{align}

El diodo parará de conducir cuando $V_m \sin \omega t$ iguale nuevamente a $V_{D0}$, permaneciendo apagado hasta el siguiente semiciclo positivo, y así sucesivamente (figura 4).

	\begin{figure}[H]
	\centering	
	\begin{tikzpicture}[scale=0.8]
		\draw (0,0) to[sV, l=${v=V_m \sin{\omega t} }$, i=$i$] ++(0,5)
		to[D] ++(2,0)
		to[battery2, invert, l=$V_{D0}$] ++(2,0)
		to[R, l=$r_D$] ++(2,0)
		to[short] ++(0.5,0)
		to[R, l_=$R_L$, v^=$v_O$, f>_=$i_L$] ++(0,-5)
		to[short] (0,0);
		
		\pgfplotsset{width=9cm,compat=1.18}
		\begin{axis}[xshift=9cm,
			yshift=-0.5,
			clip=false,
			axis y line=left,
			axis x line=middle, 
			xmin=0,xmax=2.2*pi,
			ymin=-1.5,ymax=1.5,
			xlabel=$\omega t$,
			xtick={0.3,3.14,6.28},
			xticklabels={$\alpha$, $\pi\,$, $\,\,\,2\pi$},
			ytick={0.3, 1},
			yticklabels={$V_{D0}$, $V_m$},
			]
			\addplot[smooth, red, mark=none, domain=0:2*pi, samples=200, dashed] {sin(deg(x))};
			\addplot[smooth, magenta, mark=none, domain=0.3:(pi - 0.3), samples=200] {sin(deg(x)) - 0.3};
			\addplot[color = black, dashed, thick] coordinates {(0,1) (1.57, 1)};
			\addplot[color = black, dashed, thick] coordinates {(0.3, 0) (0.3, 0.3) (0, 0.3)};
			\addplot[color = black, dashed, thick] coordinates {(2.84, 0) (2.84, 0.3) (0, 0.3)};
			\draw[->](2.6, -0.2)--(2.8, -0.05);
			\draw[<-](1.57, 1.05)--(1.8, 1.2);
			\draw[<-](2.4, 0.4)--(2.7, 0.7);
			\node[anchor=west] at (axis cs:2.1,-0.3){$\pi - \alpha$};
			\node[anchor=west] at (axis cs:1.2,1.3){$V_m \sin \omega t$};
			\node[anchor=west] at (axis cs:2.6,0.8){$v_O$};
		\end{axis}
		\draw (9.5,6) node[above right, xshift=-20pt]{$v$};
	\end{tikzpicture}
	\caption{}	
	\label{fig:figura4}
\end{figure}

Por lo tanto, el voltaje en la carga está dado por

\begin{align}
	v_O= \left\{ 
	\begin{array}{lc} \frac{R_L}{R_L + r_D} \left( V_m \sin \omega t - V_{D0} \right), & 2m\pi + \alpha \leq \omega t \leq (2m + 1)\pi - \alpha \\
		0, & otherwise
	\end{array} \right.
\end{align}

Donde $\alpha = \sin ^{-1} (V_{D0}/V_m)$ y $m=0,1, 2, ...$ \\

El ángulo $\alpha$ en el cual el diodo comienza a conducir puede hallarse con la siguiente condición

\begin{align}
	V_m \sin \alpha = V_{D0} \nonumber \\
	\sin \alpha = \frac{V_{D0}}{V_m} \nonumber
\end{align}

o

\begin{align}
	\alpha = \sin ^{-1} \left( \frac{V_{D0}}{V_m} \right) 
\end{align}

El voltaje de DC o voltaje promedio en la carga está dado por

\begin{align}
	V_{O(avg)} &= \frac{1}{2\pi} \int_{\alpha}^{\pi - \alpha} v_L d(\omega t) = \frac{1}{2\pi} \int_{\alpha}^{\pi - \alpha} \frac{R_L}{R_L + r_D} \left( V_m \sin \omega t - V_{D0} \right) d(\omega t) \nonumber \\
	&= \frac{R_L}{2\pi (R_L + r_D)} \left[ \int_{\alpha}^{\pi - \alpha} V_m \sin \omega t \, d(\omega t) - \int_{\alpha}^{\pi - \alpha} V_{D0} d(\omega t) \right] \nonumber \\
	&= \frac{R_L}{2\pi (R_L + r_D)} \left[  \left.{\left( -V_m \cos \omega t \right)}\right|_{\;\alpha}^{\;\pi - \alpha} - \left.{\left( V_{D0}(\omega t) \right)}\right|_{\;\alpha}^{\;\pi - \alpha} \right] \nonumber \\
	&= \frac{R_L}{2\pi (R_L + r_D)} \left[ -V_m \cos (\pi - \alpha) + V_m \cos \alpha - V_{D0}(\pi - \alpha) + \alpha V_{D0} \right] \nonumber \\
	&= \frac{R_L}{2\pi (R_L + r_D)} \left[ -V_m (-\cos \alpha) + V_m \cos \alpha - \pi V_{D0} + \alpha V_{DO} + \alpha V_{DO} \right] \nonumber
\end{align}

Por lo tanto

\begin{align}
	V_{O(avg)} = \frac{R_L}{2\pi (R_L + r_D)} \left[ 2V_m \cos \alpha + (2\alpha - \pi)V_{D0} \right]
\end{align}

Cuando $V_{D0} \ll V_m$ y $r_D \ll R_L$, entonces $V_{D0}$ y $r_D$ pueden ser despreciados y $v_O$ puede ser reducido a la ecuación (7). Si se llegase a invertir la conexión del diodo, la conducción ocurriría en los semiciclos negativos lo cual invertiría a su vez la polaridads de VDC.

\section*{Rectificador de onda completa}
Un puente rectificador de onda completa monofásico requiere de 4 diodos como se muestra en la figura 5. Normalmente se suele utilizar un transformador en la entrada para satisfacer los requerimientos del voltaje de salida. Dado que $v_s$ es positivo desde $\omega t = 0$ a $\pi$ y negativo desde $\omega t = \pi$ a $2\pi$, la operación del circuito se puede dividir entonces en dos intervalos: intervalo 1 e intervalo 2.

El intervalo 1 es el intervalo $0 \leq \omega t \leq \pi$ durante el semiciclo positivo del voltaje de entrada $v_s$. En este intervalo los diodos $D_3$ y  $D_4$ se encuentran polarizados en inversa. Asumiendo diodos ideales, los diodos $D_1$ y $D_2$ conducen y se comportan como corto cortocircuito. El voltaje de entrada $v_S = V_m \sin \omega t$ aparece a través de la resistencia de carga $R$. Esto es, el voltaje de salida sería:

\begin{align}
	v_O = V_m \sin \omega t, \quad 0 \leq \omega t \leq \pi
\end{align}  

El intervalo 2 es el intervalo $\pi \leq \omega t \leq 2\pi$ durante el semiciclo negativo del voltaje de entrada $v_S$. En esta ocasión los diodos $D_1$ y $D_2$ se encontrarán polarizados en inversa mientras que los diodos $D_3$ y $D_4$ conducirán, comportándose como cortocircuitos. El voltaje negativo de $v_S = V_m \sin \omega t$ aparecerá sobre la carga, es decir

\begin{align}
	v_O = -V_m \sin \omega t, \quad \pi \leq \omega t \leq 2\pi
\end{align}


	\begin{figure}[H]
	\centering	
	\begin{tikzpicture}[scale=0.8, circuitikz/diodes/scale=0.5, circuitikz/straight=true]

		\draw (0,0) to[sV, l=${v=V_m \sin{\omega t} }$] ++(0,2.62)
			  to[short] ++(3,0) node[transformer core, anchor=A1](T){}
			  (T.A2) to[short] (0,0);
	    \draw (T.B1) to [open, v=$v_s$] (T.B2);
		\draw (T.B1) to[short] ++(2,0) coordinate (top bridge);
		\draw (top bridge) to [D*,  *-*] ++(1.2,-1.32) coordinate(right bridge)
			  to[D*, *-*, invert] ++(-1.2,-1.32) coordinate (bottom bridge)
			  to[D*, *-*, invert] ++(-1.2,1.32) coordinate (left bridge)
			  to[D*, ] (top bridge);
		\draw (T.B2) to[short] (bottom bridge);
		\draw (right bridge) to[short, i=$i_O$] ++(2,0) coordinate (r1);
		\draw (left bridge) to[short] ++(-0.5,0)
			  to[short] ++(0,-3)
			  to[short] ++(4.9,0) coordinate (r2);
	    \draw (r1) to[R, v^=$v_O$, l_=$R$] (r2);
	    \node[anchor=west] at (5.9,2){$D_3$};
	    \node[anchor=west] at (8.5,2){$D_1$};
	    \node[anchor=west] at (8.5,0.5){$D_4$};
	    \node[anchor=west] at (5.9,0.5){$D_2$};
	    %\draw (p1) node[above right, xshift=-20pt]{$V_{1} = \qty{7.61}{\volt}$};
	\end{tikzpicture}
	\caption{}	
	\label{fig:figura5}
\end{figure}


	\begin{figure}[H]
	\centering	
	\begin{tikzpicture}[scale=0.8]

		\pgfplotsset{width=9cm,compat=1.18}
		\begin{axis}[
			clip=false,
			axis y line=left,
			axis x line=middle, 
			xmin=0,xmax=2.2*pi,
			ymin=-1.5,ymax=1.5,
			xlabel=$\omega t$,
			xtick={3.14,6.28},
			xticklabels={ $\pi\,$, $\,\,\,2\pi$},
			ytick={-1,1},
			yticklabels={$-V_m$, $V_m$},
			]
			\addplot[smooth, magenta, mark=none, domain=0:2*pi, samples=200] {sin(deg(x))};
			\addplot[color = black, dashed, thick] coordinates {(0,1) (1.57, 1)};
			\addplot[color = black, dashed, thick] coordinates {(0,-1) (4.8, -1)};
			\node[anchor=east] at (axis cs:0,1.5){$v_S$};
			\node[anchor=west] at (axis cs:3,0.5){$v_S = v_s = V_m \sin \omega t$};
		\end{axis}
		
		\begin{axis}[
			yshift=-7cm,
			clip=false,
			axis y line=left,
			axis x line=middle, 
			xmin=0,xmax=2.2*pi,
			ymin=-1.5,ymax=1.5,
			xlabel=$\omega t$,
			xtick={3.14,6.28},
			xticklabels={ $\pi\,$, $\,\,\,2\pi$},
			ytick={1},
			yticklabels={$V_m$},
			]
			\addplot[smooth, magenta, mark=none, domain=0:pi, samples=200] {sin(deg(x))};
			\addplot[smooth, magenta, mark=none, domain=0:pi, samples=200, xshift=3.35cm] {sin(deg(x))};
			\addplot[color = black, dashed, thick] coordinates {(0,1) (4.8, 1)};
			\addplot[color = black, dashed, thick] coordinates {(3.14,0) (3.14, 3.4)};
			\node[anchor=east] at (axis cs:0,1.5){$v_O$};
			
		\end{axis}

	\end{tikzpicture}
	\caption{}	
	\label{fig:figura6}
\end{figure}

Si $V_{D0}$ es tomado en cuenta, el voltaje pico en la carga será de $V_m - 2V_{D0}$ debido a los dos diodos en serie conduciendo.
Asdumiendo diodos ideales, el voltaje de salida promedio sería de:

\begin{align}
	V_{O(av)} &= \frac{1}{\pi} \int_{0}^{\pi} V_m \sin \omega t \, d(\omega t) \nonumber \\
	&= - \left.{\left[ \frac{V_m}{\pi} \cos({\omega t}) \right]}\right|_{\;0}^{\;\pi} \nonumber \\
	&= - \frac{V_m}{\pi} \left[ \cos \pi - \cos 0 \right] \nonumber
\end{align}

Por lo tanto:

\begin{align}
	V_{O(av)} = \frac{2V_m}{\pi} \approx 0.636Vm
\end{align}

Por lo tanto, la corriente promedio $I_{O(av)}$ en la resistencia de carga $R$ puede encontrarse mediante la ley de Ohm

\begin{align}
	I_{O(av)} = \frac{V_{O(av)}}{R} = \frac{2V_m}{\pi R} = \frac{0.636V_m}{R}
\end{align}

De la ecuación (9), el voltaje rms de una señal está dado por:

\begin{align}
	V_{rms} &= \sqrt{\frac{1}{\pi} \int_{0}^{T} { V^2_m \sin^2                                                                                                                                                                                                                                                                                                                                                                                                                                                                                                                                                                                                                                                                                                                                                                                                                                                                                                                                                                                                                                       (\omega t) d(\omega t) }} \nonumber \\
	&= \sqrt{\frac{V^2_m}{\pi} \left(\frac{\pi}{2} \right)} \nonumber 
\end{align}

\begin{align}
	V_{rms} = \frac{V_m}{\sqrt{2}} = 0.707V_m
\end{align}

La corriente de carga $I_{O(rms)}$ está dada por
\begin{align}
	I_{O(rms)} = \frac{V_{O(rms)}}{R} = \frac{V_m}{2R}
\end{align}

\subsection*{Rectificador de onda completa con el diodo real}
Nuevamente, asumiendo una aproximación lineal para los diodos en dirección directa, éste comenzara a conducir cuando $V_m \sin \omega t = V_{D0}$, la corriente de carga es de:

\begin{align}
	i_L = \frac{Vm \sin \omega t - 2V_{D0}}{R_L + 2r_D}
\end{align}

El voltaje en la carga está dado por

\begin{align}
	v_O= \left\{ 
	\begin{array}{lc} \frac{R_L}{R_L + 2r_D} \left( V_m \sin \omega t - 2V_{D0} \right), & 2m\pi + \alpha \leq \omega t \leq (2m + 1)\pi - \alpha \\
	\end{array} \right.
\end{align}

	\begin{figure}[H]
	\centering	
	\begin{tikzpicture}[scale=0.8]
		\draw (0,0) to[sV, l=${v=V_m \sin{\omega t} }$, i=$i$] ++(0,5)
		to[D] ++(2,0)
		to[battery2, invert, l=$2V_{D0}$] ++(2,0)
		to[R, l=$2r_D$] ++(2,0)
		to[short] ++(0.5,0)
		to[R, l_=$R_L$, v^=$v_O$, f>_=$i_L$] ++(0,-5)
		to[short] (0,0);
		
		\pgfplotsset{width=9cm,compat=1.18}
		\begin{axis}[xshift=9cm,
			yshift=-0.5,
			clip=false,
			axis y line=left,
			axis x line=middle, 
			xmin=0,xmax=2.2*pi,
			ymin=-1.5,ymax=1.5,
			xlabel=$\omega t$,
			xtick={0.3,3.14,6.28},
			xticklabels={$\alpha$, $\pi\,$, $\,\,\,2\pi$},
			ytick={0.3, 1},
			yticklabels={$V_{D0}$, $V_m$},
			]
			\addplot[smooth, red, mark=none, domain=0:2*pi, samples=200, dashed] {sin(deg(x))};
			\addplot[smooth, green, mark=none, domain=0:pi, samples=200] {sin(deg(x))};
			\addplot[smooth, green, mark=none, domain=0:pi, samples=200, xshift=3.35cm] {sin(deg(x))};
			\addplot[smooth, magenta, mark=none, domain=0.3:(pi - 0.3), samples=200] {sin(deg(x)) - 0.3};
			\addplot[smooth, magenta, mark=none, domain=0.3:(pi - 0.3), samples=200, xshift=3.35cm] {sin(deg(x)) - 0.3};
			\addplot[color = black, dashed, thick] coordinates {(0,1) (1.57, 1)};
			\addplot[color = black, dashed, thick] coordinates {(0.3, 0) (0.3, 0.3) (0, 0.3)};
			\addplot[color = black, dashed, thick] coordinates {(2.84, 0) (2.84, 0.3) (0, 0.3)};
			\draw[->](2.6, -0.2)--(2.8, -0.05);
			\draw[<-](1.57, 1.05)--(1.8, 1.2);
			\draw[<-](2.4, 0.4)--(2.7, 0.7);
			\draw[<-](4.6, 1.05)--(4.8, 1.2);
			\node[anchor=west] at (axis cs:2.1,-0.3){$\pi - \alpha$};
			\node[anchor=west] at (axis cs:1.2,1.3){$V_m \sin \omega t$};
			\node[anchor=west] at (axis cs:4.2,1.3){$V_O \, ideal$};
			\node[anchor=west] at (axis cs:2.6,0.8){$v_O$};
		\end{axis}
		\draw (9.5,6) node[above right, xshift=-20pt]{$v$};
	\end{tikzpicture}
	\caption{}	
	\label{fig:figura7}
\end{figure}

El voltaje de DC estará dado por:

\begin{align}
	V_{O(avg)} &= \frac{1}{\pi} \int_{\alpha}^{\pi - \alpha} v_L d(\omega t) = \frac{1}{\pi} \int_{\alpha}^{\pi - \alpha} \frac{R_L}{R_L + 2r_D} \left( V_m \sin \omega t - 2V_{D0} \right) d(\omega t) \nonumber \\
	&= \frac{R_L}{\pi (R_L + 2r_D)} \left[ \int_{\alpha}^{\pi - \alpha} V_m \sin \omega t \, d(\omega t) - \int_{\alpha}^{\pi - \alpha} 2V_{D0} d(\omega t) \right] \nonumber \\
	&= \frac{R_L}{\pi (R_L + 2r_D)} \left[  \left.{\left( -V_m \cos \omega t \right)}\right|_{\;\alpha}^{\;\pi - \alpha} - \left.{\left( 2V_{D0}(\omega t) \right)}\right|_{\;\alpha}^{\;\pi - \alpha} \right] \nonumber \\
	&= \frac{R_L}{\pi (R_L + 2r_D)} \left[ -V_m \cos (\pi - \alpha) + V_m \cos \alpha - 2V_{D0}(\pi - \alpha) + 2\alpha V_{D0} \right] \nonumber \\
	&= \frac{R_L}{\pi (R_L + 2r_D)} \left[ -V_m (-\cos \alpha) + V_m \cos \alpha - 2\pi V_{D0} + 2\alpha V_{DO} + 2\alpha V_{DO} \right] \nonumber
\end{align}

Por lo tanto

\begin{align}
	V_{O(avg)} = \frac{R_L}{\pi (R_L + r_D)} \left[ 2V_m \cos \alpha + 2(2\alpha - \pi)V_{D0} \right]
\end{align}


\end{document}